\bookchapter{제7장: 최첨단 기능 완전 해설}

\chapterintro{기술은 매일 발전하고, 스마트폰 기능도 계속 진화합니다. 위성 통신, 무통신, 차량 사고 감지, SparkLink 기술... 이처럼 과학 소설처럼 들리는 기능들이 이미 우리 생활에 들어와 있습니다. 이 최첨단 기술들을 함께 알아봅시다.}

\booksection{위성 통신: 신호가 없어도 구조 요청}

위성 통신은 최근 스마트폰 분야의 큰 돌파구입니다. 황야, 깊은 산림에서 스마트폰 신호가 전혀 없을 때, 위성 통신이 구명줄이 될 수 있습니다.

\booksubsection{작동 원리:}

스마트폰이 궤도상의 통신 위성과 직접 연결하여 구조 메시지를 전송합니다. 현재는 주로 문자와 위치 정보 전송만 지원하며, 전화 통화는 불가능합니다.

\booksubsection{지원 기종:}

\begin{itemize}
  \item \textbf{Apple}: iPhone 14 시리즈 이후, 위성을 통한 긴급 SOS 지원. iPhone 17 시리즈는 위성을 통한 iMessage 및 문자 전송을 추가로 지원하며, 더 많은 지역의 커버리지 확장.
  \item \textbf{화웨이}: Mate 50 시리즈부터 베이더우 위성 메시지(문자만) 지원, Mate 60 Pro 시리즈는 텐통 위성 통화(전화 가능) 지원, Mate 70 시리즈는 위성 통신 능력이 더욱 강화됨.
  \item \textbf{샤오미/OPPO/vivo}: 2024년부터 플래그십 기기에 베이더우 위성 메시지 기능 점차 지원.
  \item \textbf{발전 추세}: 위성 통신이 고급 플래그십 전용 기능에서 점차 보급되어, 미래에는 중급 기기도 지원할 전망. 중국 베이더우 위성 시스템의 지원이 국산 스마트폰에서 가장 광범위.
\end{itemize}

\booksubsection{사용 시나리오:}

\begin{itemize}
  \item 아웃도어 탐험, 등산 중 구조 요청.
  \item 해상 조난.
  \item 자연 재해로 지상 통신 두절.
\end{itemize}

\booksubsection{사용 방법:}

\begin{enumerate}
  \item 반드시 탁 트인 곳에서, 하늘을 볼 수 있어야 합니다.
  \item 스마트폰의 위성 통신 기능을 켜세요(일반적으로 긴급 통화 인터페이스에서).
  \item 화면 안내에 따라, 위성 방향으로 맞추세요.
  \item 프리셋된 구조 요청 정보와 위치를 보내기.
\end{enumerate}

\begin{tcbinfo}{작은 지식}
위성 통신은 스마트폰에 특수 하드웨어 안테나가 필요하며, 모든 스마트폰이 지원하는 것은 아닙니다. 또한, 위성 통신은 일반적으로 유료 구독 서비스가 필요하지만, 긴급 구조 요청 기능은 일반적으로 무료입니다.

\end{tcbinfo}

\booksection{무선망 통신: 근거리 통신 새로운 방식}

“무선망 통신”은 이동 네트워크와 Wi-Fi가 없는 상황에서, 두 대의 스마트폰 사이에 직접 통신하는 기술을 말합니다.

\booksubsection{기술 원리:}

Bluetooth 또는 전용 무선 주파수를 이용하여, 단거리(일반적으로 수백 미터에서 수 킬로미터)내에 피어투피어 연결을 설정합니다.

\booksubsection{앱 시나리오:}

\begin{itemize}
  \item \textbf{야외 활동}:트레킹,캠핑 시 팀원 간 통신.
  \item \textbf{재난 구조}:지진,홍수 후 통신 단절 시.
  \item \textbf{대형 활동}:콘서트,체육 경기장은 사람이 많아 신호가 나쁠 시.
\end{itemize}

\booksubsection{지원 플랫폼:}

\begin{itemize}
  \item \textbf{애플}:iOS 17의 "NameDrop"과 "오프라인 친구 찾기" 기능.
  \item \textbf{화웨이}:창련 App이 무선망 통화를 지원.
  \item \textbf{샤오미}:샤오미 묘향 센터의 일부 기능.
\end{itemize}

\booksection{차량 사고 감지:위급 순간 목숨 구하기}

차량 사고 감지는 스마트폰 센서를 이용해 심각한 차량 사고를 감지하고, 자동으로 긴급 전화를 거는 기능입니다.

\booksubsection{작동 원리:}

스마트폰은 가속도계,자이로스코프,기압계,마이크 등 여러 센서를 통해, 심각한 충돌이 발생했는지 종합적으로 판단합니다. 차량 사고를 감지하면, 스마트폰은:

\begin{enumerate}
  \item 카운트다운 인터페이스를 표시하여, 괜찮은지 묻습니다.
  \item 만약 카운트다운이 끝날 때까지 응답이 없으면, 자동으로 긴급 전화를 겁니다.
  \item 긴급 연락처에 당신의 위치를 보냅니다.
\end{enumerate}

\booksubsection{지원 기종:}

\begin{itemize}
  \item \textbf{애플}:iPhone 14 시리즈 및 이후.
  \item \textbf{안드로이드}:일부 고급 기종이 순차적으로 지원.
\end{itemize}

\booksubsection{주의 사항:}

\begin{itemize}
  \item 기능을 켜야 사용할 수 있습니다.
  \item 오탐지될 수 있습니다(예: 롤러코스터,격렬한 운동), 제때 취소하면 됩니다.
  \item 센서와 알고리즘에 의존하여 100\% 신뢰할 수는 없지만, 한 겹의 보장이 더해집니다.
\end{itemize}

\booksection{SparkLink(SparkLink) 기술: 차세대 무선 연결}

SparkLink(SparkLink)는 화웨이가 주도하여 연구 개발한 차세대 무선 통신 기술로, 목표는 일부 시나리오에서 Bluetooth와 Wi-Fi를 대체하는 것입니다. 핵심 지표에서 기존 방안에 대응하거나 심지어 능가합니다: 공중 인터페이스 속도는 900Mbps에 달하고(Bluetooth의 50Mbps를 훨씬 상회), 지연은 20 $\mu$ s까지 낮으며(Bluetooth 약 15ms), 단일 네트워크는 4096개 기기의 동시 연결을 지원합니다(Bluetooth는 일반적으로 8개 이내).

\booksubsection{기술 특징:}

\begin{itemize}
  \item \textbf{초저지연}:종단 간 지연은 20마이크로초(0.02ms)로, Bluetooth보다 두 자릿수 낮아, 게임과 음악 동기화를 거의 인지하지 못합니다.
  \item \textbf{고속}: 단일 링크 피크 속도 900Mbps, 블루투스 오디오의 수십 배에 달하며, 무손실 오디오 또는 HD 영상 전송 가능.
  \item \textbf{저전력}: 대기 전력 소모는 블루투스와 비슷하지만, 전송 효율이 더 높아 단위 데이터당 에너지 소모가 더 낮음.
  \item \textbf{광범위 커버리지+고동시성}:연결 거리는 Bluetooth보다 50\% 이상 길고, 단일 네트워크는 수천 기기의 동시 온라인을 지원합니다.
\end{itemize}

\booksubsection{앱 시나리오:}

\begin{itemize}
  \item \textbf{무선 이어폰}:더 낮은 지연, 게임 시 음향-화면 동기화.
  \item \textbf{스마트워치}: 더 안정적인 연결.
  \item \textbf{스마트 홈}: 더 많은 기기 동시 연결.
  \item \textbf{차량 연동}: 더 매끄러운 스마트폰-차량 연결.
\end{itemize}

\booksubsection{지원기기:}

현재 주로 화웨이 생태계 기기입니다, 예를 들어 화웨이 Mate 60 시리즈,일부 화웨이 이어폰과 워치. 미래에는 더 많은 브랜드가 지원할 것입니다.

\textbf{기술 전망}\par

SparkLink 기술은 무선 통신의 미래 방향을 대표합니다. 더 많은 제조사가 가입함에 따라, 우리는 더 많은 SparkLink 지원 기기를 보게 되어, 더 나은 무선 경험을 누릴 것입니다.

\begin{tcbpractice}{한번 확인해보세요: 당신의 스마트폰이 이러한 기능을 지원하는지}
\begin{enumerate}
  \item 설정에서 "위성"을 검색하여, 위성 통신 관련 옵션이 있는지 확인하세요.
  \item "차량 사고 감지" 또는 "긴급 SOS"를 검색하여, 켜져 있는지 확인하세요.
  \item 만약 당신의 스마트폰이 화웨이라면, SparkLink 기기 연결을 지원하는지 확인하세요.
  \item 이러한 기능의 사용 방법을 알아두면, 위급 순간에 도움이 될 수 있습니다.
\end{enumerate}

\end{tcbpractice}
